% Section 5.4, from lectures/L05/README.md: the objectives, the evaluation questions, and the next
% lecture.

\section{Review}
\label{c5:sec:review}

\needspace{6\baselineskip}
\subsection*{What you should be able to do now}
After this chapter, you should:
\begin{itemize}
  \item Understand what templates are and why they are used.
  \item Understand how function templates enable generic code.
  \item Understand how class templates can be used for reusable components (e.g.\ buffers,
    frames).
  \item Understand that templates are instantiated at compile time and generate concrete
    types and functions.
  \item Understand how templates affect code size in embedded systems.
  \item Understand when templates are appropriate in embedded systems and when they are not.
  \item Be able to read and follow simple template-based code.
  \item Understand the purpose of type traits for basic compile-time checks.
\end{itemize}

\needspace{6\baselineskip}
\subsection*{Questions to test yourself}
\begin{enumerate}
  \item What problem do templates solve in C++?
  \item What is the difference between a function template and a class template?
  \item What does it mean that templates are instantiated at compile time, and what does the
    compiler generate?
  \item Why can templates increase code size in embedded systems?
  \item Give an example where a template is more suitable than an interface.
  \item Why must template implementations typically be placed in header files?
\end{enumerate}

\needspace{6\baselineskip}
\subsection*{Looking ahead}
Chapter~6 is about concurrency concepts in modern C++:
\begin{itemize}
  \item Threads.
  \item Mutexes.
  \item Atomic variables.
  \item Lock guards.
  \item Synchronization primitives.
  \item Thread-safe design.
\end{itemize}
